% GENERATED FILE. Edit canonical lessons, then run npm run generate:books.
% canonical-source-hash: fnv1a64:a0df70fa4e2cc80e
% canonical-lessons: PA-C268-adda, PA-C268-kandaktar, PA-C268-sit, PA-C268-bhir, PA-C268-tikna

\chapter{At the Bus Stand}
\label{ch:pa-a2-at-the-bus-stand}
\hlchaptermodality{268}

\begin{chapteropening}
\textbf{By the end of this chapter:} \emph{I can find a seat at a busy bus stand.}

Complete the last lesson of chapter 268: I can find a seat at a busy bus stand.
\end{chapteropening}

\section[\texorpdfstring{\pa{ਅੱਡਾ} (aḍḍā)}{ਅੱਡਾ (aḍḍā)}]{\pa{ਅੱਡਾ} (aḍḍā) — a bus stand; a stand for vehicles}

\label{lesson:PA-C268-adda}



\begin{quote}
Before the new one: say the Punjabi for a court of law, then the Punjabi for whereabouts.
\end{quote}

\subsection*{You'll want to know: \pa{ਅੱਡਾ}}
\textbf{\pa{ਅੱਡਾ}} — \emph{aḍḍā} — \textquotedblleft{}a bus stand; a stand for vehicles\textquotedblright{}.

\begin{quote}
\pa{ਬੱਸ} \pa{ਅੱਡਾ} \pa{ਸ਼ਹਿਰ} \pa{ਦੇ} \pa{ਬਾਹਰ} \pa{ਹੈ।}
\end{quote}

\begin{quote}
\emph{bass aḍḍā shahir de bāhar hai.}
\end{quote}

\begin{quote}
The bus stand is outside the town.
\end{quote}

A masculine noun.

\begin{cousinweb}
\textbf{\pa{ਅੱਡਾ}} (\emph{aḍḍā}) is inherited from a reconstructed \emph{*aḍḍ}, \textquotedblleft{}obstruct, stop\textquotedblright{}; CDIAL lists Punjabi \emph{aḍḍā}, \textquotedblleft{}perch for birds\textquotedblright{}, under it (CDIAL 188). The sense \textquotedblleft{}a stand\textquotedblright{} is not given there.
\end{cousinweb}

\begin{grammarlens}[title={What you've built}]
One more word for this chapter.
\end{grammarlens}

\subsection*{Your turn}
\begin{itemize}
\raggedright
  \item \emph{Say it:} \emph{aḍḍā}
  \item \emph{Say it:} \emph{aḍḍā}, once more
  \item \emph{Say it:} \emph{bass aḍḍā shahir de bāhar hai.} — The bus stand is outside the town.
  \item \emph{Recall:} say the Punjabi for a court of law, then the Punjabi for whereabouts, then say \emph{aḍḍā} again
\end{itemize}

\begin{tcolorbox}[breakable,colback=teal!4,colframe=teal!35!black,title={Before you move on}]
What does \pa{ਅੱਡਾ} mean? (\textquotedblleft{}A bus stand; a stand for vehicles\textquotedblright{}.) Say it once more.
\end{tcolorbox}
\section[\texorpdfstring{\pa{ਕੰਡਕਟਰ} (kaṇḍakṭar)}{ਕੰਡਕਟਰ (kaṇḍakṭar)}]{\pa{ਕੰਡਕਟਰ} (kaṇḍakṭar) — a bus conductor}

\label{lesson:PA-C268-kandaktar}



\begin{quote}
Before the new one: say the Punjabi for whereabouts, then the Punjabi for a bus stand.
\end{quote}

\subsection*{You'll want to know: \pa{ਕੰਡਕਟਰ}}
\textbf{\pa{ਕੰਡਕਟਰ}} — \emph{kaṇḍakṭar} — \textquotedblleft{}a bus conductor\textquotedblright{}.

\begin{quote}
\pa{ਕੰਡਕਟਰ} \pa{ਬੱਸ} \pa{ਵਿੱਚ} \pa{ਟਿਕਟਾਂ} \pa{ਵੇਚਦਾ} \pa{ਹੈ।}
\end{quote}

\begin{quote}
\emph{kaṇḍakṭar bass vich ṭikṭāṁ vechdā hai.}
\end{quote}

\begin{quote}
The conductor sells tickets on the bus.
\end{quote}

A masculine noun.

\begin{cousinweb}
\textbf{\pa{ਕੰਡਕਟਰ}} (\emph{kaṇḍakṭar}) is borrowed from English \emph{conductor}.
\end{cousinweb}

\begin{grammarlens}[title={What you've built}]
One more word for this chapter.
\end{grammarlens}

\subsection*{Your turn}
\begin{itemize}
\raggedright
  \item \emph{Say it:} \emph{kaṇḍakṭar}
  \item \emph{Say it:} \emph{kaṇḍakṭar}, once more
  \item \emph{Say it:} \emph{kaṇḍakṭar bass vich ṭikṭāṁ vechdā hai.} — The conductor sells tickets on the bus.
  \item \emph{Recall:} say the Punjabi for whereabouts, then the Punjabi for a bus stand, then say \emph{kaṇḍakṭar} again
  \item \emph{Write it:} copy \textbf{\pa{ਕੰਡਕਟਰ}} once, checking the whole word
\end{itemize}

\begin{tcolorbox}[breakable,colback=teal!4,colframe=teal!35!black,title={Before you move on}]
What does \pa{ਕੰਡਕਟਰ} mean? (\textquotedblleft{}A bus conductor\textquotedblright{}.) Say it once more.
\end{tcolorbox}
\section[\texorpdfstring{\pa{ਸੀਟ} (sīṭ)}{ਸੀਟ (sīṭ)}]{\pa{ਸੀਟ} (sīṭ) — a seat}

\label{lesson:PA-C268-sit}



\begin{quote}
Before the new one: say the Punjabi for a bus stand, then the Punjabi for a bus conductor.
\end{quote}

\subsection*{You'll want to know: \pa{ਸੀਟ}}
\textbf{\pa{ਸੀਟ}} — \emph{sīṭ} — \textquotedblleft{}a seat\textquotedblright{}.

\begin{quote}
\pa{ਕੀ} \pa{ਬੱਸ} \pa{ਵਿੱਚ} \pa{ਸੀਟ} \pa{ਖ਼ਾਲੀ} \pa{ਹੈ}?
\end{quote}

\begin{quote}
\emph{kī bass vich sīṭ khālī hai?}
\end{quote}

\begin{quote}
Is there a free seat on the bus?
\end{quote}

A feminine noun.

\begin{cousinweb}
\textbf{\pa{ਸੀਟ}} (\emph{sīṭ}) is borrowed from English \emph{seat}.
\end{cousinweb}

\begin{grammarlens}[title={What you've built}]
One more word for this chapter.
\end{grammarlens}

\subsection*{Your turn}
\begin{itemize}
\raggedright
  \item \emph{Say it:} \emph{sīṭ}
  \item \emph{Say it:} \emph{sīṭ}, once more
  \item \emph{Say it:} \emph{kī bass vich sīṭ khālī hai?} — Is there a free seat on the bus?
  \item \emph{Recall:} say the Punjabi for a bus stand, then the Punjabi for a bus conductor, then say \emph{sīṭ} again
  \item \emph{Write it:} copy \textbf{\pa{ਸੀਟ}} once, checking the whole word
\end{itemize}

\begin{tcolorbox}[breakable,colback=teal!4,colframe=teal!35!black,title={Before you move on}]
What does \pa{ਸੀਟ} mean? (\textquotedblleft{}A seat\textquotedblright{}.) Say it once more.
\end{tcolorbox}
\section[\texorpdfstring{\pa{ਭੀੜ} (bhīṛ)}{ਭੀੜ (bhīṛ)}]{\pa{ਭੀੜ} (bhīṛ) — a crowd}

\label{lesson:PA-C268-bhir}



\begin{quote}
Before the new one: say the Punjabi for a bus conductor, then the Punjabi for a seat.
\end{quote}

\subsection*{You'll want to know: \pa{ਭੀੜ}}
\textbf{\pa{ਭੀੜ}} — \emph{bhīṛ} — \textquotedblleft{}a crowd\textquotedblright{}.

\begin{quote}
\pa{ਸ਼ਾਮ} \pa{ਨੂੰ} \pa{ਅੱਡੇ} \pa{ਉੱਤੇ} \pa{ਬਹੁਤ} \pa{ਭੀੜ} \pa{ਹੁੰਦੀ} \pa{ਹੈ।}
\end{quote}

\begin{quote}
\emph{shām nū̃ aḍḍe utte bahut bhīṛ hundī hai.}
\end{quote}

\begin{quote}
In the evening there is a big crowd at the stand.
\end{quote}

A feminine noun.

\begin{cousinweb}
\textbf{\pa{ਭੀੜ}} (\emph{bhīṛ}) is inherited from a reconstructed \emph{*bhiṭ}, \textquotedblleft{}meet, throng\textquotedblright{}; CDIAL lists Punjabi \emph{bhīṛ}, \textquotedblleft{}crowd\textquotedblright{}, under it (CDIAL 9490).
\end{cousinweb}

\begin{grammarlens}[title={What you've built}]
One more word for this chapter.
\end{grammarlens}

\subsection*{Your turn}
\begin{itemize}
\raggedright
  \item \emph{Say it:} \emph{bhīṛ}
  \item \emph{Say it:} \emph{bhīṛ}, once more
  \item \emph{Say it:} \emph{shām nū̃ aḍḍe utte bahut bhīṛ hundī hai.} — In the evening there is a big crowd at the stand.
  \item \emph{Recall:} say the Punjabi for a bus conductor, then the Punjabi for a seat, then say \emph{bhīṛ} again
  \item \emph{Write it:} copy \textbf{\pa{ਭੀੜ}} once, checking the whole word
\end{itemize}

\begin{tcolorbox}[breakable,colback=teal!4,colframe=teal!35!black,title={Before you move on}]
What does \pa{ਭੀੜ} mean? (\textquotedblleft{}A crowd\textquotedblright{}.) Say it once more.
\end{tcolorbox}
\section[\texorpdfstring{\pa{ਟਿਕਣਾ} (ṭikṇā)}{ਟਿਕਣਾ (ṭikṇā)}]{\pa{ਟਿਕਣਾ} (ṭikṇā) — to stay put; to stay for a while}

\label{lesson:PA-C268-tikna}



\begin{quote}
Before the new one: say the Punjabi for a seat, then the Punjabi for a crowd.
\end{quote}

\subsection*{You'll want to know: \pa{ਟਿਕਣਾ}}
\textbf{\pa{ਟਿਕਣਾ}} — \emph{ṭikṇā} — \textquotedblleft{}to stay put; to stay for a while\textquotedblright{}.

\begin{quote}
\pa{ਅਸੀਂ} \pa{ਸ਼ਹਿਰ} \pa{ਵਿੱਚ} \pa{ਦੋ} \pa{ਦਿਨ} \pa{ਟਿਕਾਂਗੇ।}
\end{quote}

\begin{quote}
\emph{asīṁ shahir vich do din ṭikāṅge.}
\end{quote}

\begin{quote}
We will stay two days in the town.
\end{quote}

Past (I …): \pa{ਮੈਂ} \pa{ਅੱਡੇ} \pa{ਉੱਤੇ} \pa{ਟਿਕਿਆ}, \emph{maiṁ aḍḍe utte ṭikiā}.

\begin{cousinweb}
\textbf{\pa{ਟਿਕਣਾ}} (\emph{ṭikṇā}) is inherited from a reconstructed \emph{*ṭakk}, \textquotedblleft{}remain, stop\textquotedblright{}; CDIAL lists Punjabi \emph{ṭikṇā}, \textquotedblleft{}to stay\textquotedblright{}, under it (CDIAL 5420).
\end{cousinweb}

\begin{grammarlens}[title={What you've built}]
One more word for this chapter.
\end{grammarlens}

\subsection*{Your turn}
\begin{itemize}
\raggedright
  \item \emph{Say it:} \emph{ṭikṇā}
  \item \emph{Say it:} \emph{ṭikṇā}, once more
  \item \emph{Say it:} \emph{asīṁ shahir vich do din ṭikāṅge.} — We will stay two days in the town.
  \item \emph{Recall:} say the Punjabi for a seat, then the Punjabi for a crowd, then say \emph{ṭikṇā} again
  \item \emph{Write it:} copy \textbf{\pa{ਟਿਕਣਾ}} once, checking the whole word
\end{itemize}

\begin{tcolorbox}[breakable,colback=teal!4,colframe=teal!35!black,title={Before you move on}]
What does \pa{ਟਿਕਣਾ} mean? (\textquotedblleft{}To stay put; to stay for a while\textquotedblright{}.) Say it once more.
\end{tcolorbox}
